\section{Waymo: ingeniería de sistemas y mentalidad de ingeniero}

Al terminar el doctorado eligió la conducción autónoma, porque en ese momento era la forma más clara de physical world AI. Lo que más aprendió en Waymo no fueron solo los algoritmos de conducción autónoma, sino cómo funciona el sistema completo: cómo se integran la percepción, la localización, los mapas de alta precisión fuera de línea, la predicción, la toma de decisiones, la planificación, el control, la simulación y la cadena de herramientas en la nube hasta formar un sistema confiable.

\subsection{La continuidad histórica de la arquitectura de conducción autónoma}

Al releer los artículos de conducción autónoma de alrededor de 2008, encontró que la arquitectura general de hacia 2018 no difería en lo fundamental de los primeros marcos: el sistema seguía dividido en módulos de percepción, localización, mapas, planificación y control. El verdadero cambio fue que, cuando las técnicas de IA maduraron, muchos módulos de percepción pasaron de clustering, reglas y métodos tradicionales a redes neuronales.

De ahí surge una distinción importante: la lógica de fondo de la gran arquitectura tradicional de conducción autónoma se parece más a la de robotics y da prioridad a la modularidad, la interpretabilidad y los corner case; la vía AI native da prioridad al enfoque basado en datos, al diseño de extremo a extremo y a la mejora del benchmark global. No se trata de decidir cuál tiene razón, sino de equilibrar de manera distinta la confiabilidad, la capacidad de iteración, la interpretabilidad y la escalabilidad.

\begin{warningbox}{Lo modular no es atraso, ni lo de extremo a extremo lo resuelve todo}
En un sistema modular es más fácil localizar los corner case y los límites de responsabilidad, pero puede tardar en cambiar de rumbo y acumular reglas locales; un sistema de extremo a extremo mejora con más facilidad las métricas globales a partir de datos, pero puede introducir fallas que no se pueden explicar. La IA del mundo físico a menudo necesita volver a trazar la frontera entre ambos.
\end{warningbox}

\subsection{Mentalidad de ingeniero: descomponer y medir}

Gao Jiyang (高继扬) resume que la formación central que le dio Waymo fue la mentalidad de ingeniero: descomponer y medir. Se divide un problema complejo en subproblemas manejables, y estos a su vez en código y pruebas; al mismo tiempo, se construye un sistema de medición que va de las métricas de nivel superior a las intermedias y hasta las pruebas unitarias de base. Esto difiere de resolver problemas en las olimpiadas de física, donde se trata sobre todo de reconocer a qué tipo de problema corresponde cada uno; un sistema de ingeniería, en cambio, tiene que funcionar de manera continua, localizar problemas de manera continua y volver de manera continua al objetivo global.

\begin{importantbox}{Mentalidad de ingeniero}
La mentalidad de ingeniero no consiste en «saber escribir código», sino en descomponer un sistema complejo en módulos verificables y conectar, mediante un sistema de métricas, los cambios locales con los resultados de nivel superior.
\end{importantbox}

\subsection{Resumen de la sección}

Waymo le dio a Gao Jiyang una visión del sistema completo y una formación en ingeniería: saber cómo se descompone un sistema de IA del mundo físico, cómo se mide y cómo se ejecuta a largo plazo. Pero estaba lejos de los clientes, del producto y de la gestión de una empresa, así que su siguiente paso fue ir a Momenta, más cerca de la entrega de producción en serie.
